% =====================================================================
%  sec3.tex — 3. The Hitchin Fibration（含 3.1 A Few Hints）
%  原文 pp.3–5
% =====================================================================
\section{Hitchin 纤维化}

我们将用到 Hitchin 纤维化。这是全纯依赖于复结构 $I$ 的映射：它把 Higgs 丛 $(E,\varphi)$ 送到 $\varphi$ 的特征多项式。例如当 $G=\SU(2)$ 时，$(E,\varphi)$ 被送到二次微分 $\det\varphi$\zhnt{3.1}。于是 Hitchin 纤维化在此情形下的目标，就是以二次微分为参数的空间 $B=H^{0}(C,K^{2})$。对任何 $G$ 这都有自然的类比\zhnt{3.2}。从复结构 $I$ 的角度看，映射 $\pi:\MH\to B$ 的一般纤维是复阿贝尔簇（说得稍准确些，一般地是某个阿贝尔簇上的 torsor）\zhnt{3.3}。从全纯辛形式 $\Omega_{I}$ 的角度看，这些纤维是拉格朗日的。这样一种以复拉格朗日环面为纤维的纤维化，使 $\MH$ 成为一个复意义下的完全可积 Hamilton 系统，见 Hitchin (1987b)。

不过，对 Hitchin 纤维化还有另一种看法，最早由 Hausel 与 Thaddeus (2002) 描述。让我们回到以实辛结构 $\omega$ 定义的 A-模型。既然 $\pi:\MH\to B$ 的纤维关于 $\Omega_{I}=\omega_{J}+i\omega_{K}$ 是拉格朗日的，它们特别地关于 $\omega=\omega_{K}$ 也是拉格朗日的。而且，这些纤维关于复结构 $I$ 是全纯的，因此实际上在同调类中是面积最小的——这里面积用 $\MH$ 上的超凯勒度量计算\zhnt{3.4}。于是，从 A-模型的角度看，Hitchin 纤维化实际上是 $\MH$ 的一个特殊拉格朗日环面纤维化。

人们相信，镜像对称来自特殊拉格朗日纤维化在纤维方向上的 T-对偶，见 Strominger 等 (1996)\zhnt{3.5}。一般而言，要显式地给出一个非平凡的特殊拉格朗日纤维化是极其困难的。目前的例子是极少数确实做得到的例子之一，靠的是 $\MH$ 的超凯勒结构及其可积性质。非平凡的特殊拉格朗日纤维化之所以难以理解，是因为难以弄清楚其奇点的结构。在超凯勒的语境中，这些纤维在另一个复结构下是全纯的，这一事实使一切都容易把握得多。

一旦真的找到了一个特殊拉格朗日纤维化，要给出镜像对称的一个例子，我们该做的就是构造对偶的特殊拉格朗日纤维化，它应当镜像于原来的纤维化。镜像映射把一侧的辛结构与另一侧的复结构互换。

在目前的语境中，对偶纤维化有一个非常漂亮的描述：正如 Hausel 与 Thaddeus (2002) 首先证明的，它恰好就是对偶群的 Hitchin 纤维化。于是我们考虑 $\MH(\LG,C)$，即对偶群 $\LG$ 的 Hitchin 方程之解的模空间。结果，$G$ 与 $\LG$ 的 Hitchin 纤维化的底可以自然地等同。由此得到的图景大致如下（图~\ref{fig:hitchin}）。

\begin{figure}[htbp]
\centering
\begin{tikzpicture}[>={Stealth[length=2.2mm]},scale=1.0]
  \node (L) at (0,1.05) {$\MH(\LG,C)$};
  \node (G) at (5.0,1.05) {$\MH(G,C)$};
  \node (B) at (2.5,0) {$B$};
  \draw[->] (L) -- (B);
  \draw[->] (G) -- (B);
\end{tikzpicture}
\caption{Hitchin 纤维化：$G$ 与对偶群 $\LG$ 的解空间以同一个底 $B$ 为底。（原文第 4 页以文字排出的示意图，此处重排为图形，内容未变。）}
\label{fig:hitchin}
\end{figure}

在复结构 $I$ 下，一般点 $b\in B$ 上的纤维粗略地说是对偶阿贝尔簇（更准确地说，是对偶阿贝尔簇上的 torsor）。换一个角度，这些纤维是辛结构 $\omega=\omega_{K}$ 下的特殊拉格朗日子流形。从后一角度看，同一幅图给出的是 $\MH(\LG,C)$ 在复结构 $J$ 下的 B-模型与 $\MH(G,C)$ 以辛结构 $\omega_{K}$ 的 A-模型之间的镜像对称。

正如我们刚刚解释的，使这一镜像对称显露出来的工具，是 $\MH$ 的超凯勒结构及其 Hitchin 纤维化。这些结构依赖于 $C$ 上复结构的选择，但事实上所得的镜像对称并不真正依赖这一选择\zhnt{3.6}。这一点是在 Kapustin 与 Witten (2007) 中、在从四维拓扑场论导出这一镜像对称例子的过程中被证明的。这里所说的拓扑场论，由 $\mathcal{N}=4$ 超 Yang–Mills 理论经扭转得到。

\subsection{几点提示}

要从这个镜像对称的例子走到几何 Langlands 对偶的通常表述，还有几个障碍要克服。遗憾的是，此处不便多谈，只能给出几点提示。

一个要点是：在通常的表述中，$\MH(\LG,C)$ 上一个 B-膜的对偶，应当不是 $\MH(G,C)$ 上的 A-膜——而那恰恰是我们从上述构造中最自然地得到的东西——而是 $M(G,C)$（$C$ 上 $G$-丛的模空间）上的 D-模层（所谓 D-模层，按定义就是 $M(G,C)$ 上微分算子层 $\mathcal{D}$ 的模层）\zhnt{3.7}。这两种陈述之间的联系在 Kapustin 与 Witten (2007) 第 11 节中得到解释，用到了 $\MH(G,C)$ 上一个特殊的 A-膜，它与 $M(G,C)$ 上的微分算子有密切关系。这种联系之所以可能，是因为如 Hitchin (1987a) 所述，$\MH(G,C)$ 以 $T^{*}M_{\mathrm{st}}(G,C)$ 作为一个 Zariski 开集；这里 $M_{\mathrm{st}}(G,C)$ 是 $M(G,C)$ 中参数化严格稳定丛的子空间。

另一个要点如下。在通常的表述中，几何 Hecke 算子起着核心作用，它们作用在 $C$ 上的全纯 $G$-丛上，因而也作用在 $M(G,C)$ 上的 D-模上。它们在这一图景中有其自然的位置，但正是在这里，若只试图用二维 $\sigma$ 模型与镜像对称来表达这一主题，就会漏掉一些东西。这一特定的镜像对称例子实际上源自一个底层四维规范场论中的对偶。一旦理解了这一点，关于规范场论的 Wilson 与 't Hooft 线算子的基本事实就会给出关于 Hecke 特征层的通常陈述，如 Kapustin 与 Witten (2007) 第 9、10 节所解释的那样。几何 Hecke 算子在这一语境中自然地用三维规范场论的 Bogomolny 方程重新解释；这些方程在几何上与物理上都极有意义，并已被大量研究，例如见 Atiyah 与 Hitchin (1988)\zhnt{3.8}。

这些陈述中的某一些若要得到恰当的表述，就会引出四维的另一个重要作用。几何 Langlands 的通常表述所涉及的 D-模，不是在半稳定 $G$-丛的模空间上，而是在所有 $G$-丛的模叠上。主要原因是：若只考虑半稳定丛，就看不到 Hecke 算子的作用。正如我们将在第 6 节说明的，叠在标准描述中的这种地位，可以被理解为一条强有力的线索，指向一个从四维规范场论出发的替代进路。
